Kubernetes has developed as the standard for workload orchestration.
However, with Ansible being a popular tool to manage stateful workloads such as virtual machines, a bridge is needed to combine both under a common access plane.
Because Ansible is imperative and imposition based, while Kubernetes operates on a promise based model, combining the two systems under the Kubernetes API requires making a promise to impose.
After defining requirements, a Kubernetes operator was designed and implemented that allows scheduling Ansible executions and defining an Ansible inventory from the Kubernetes API.
The notable exposed CRDs are \anhost and \angroup used to declare the inventory, \anplaybook pointing the operator at the playbook source, either inlined into the resource or pulled from git, as well as \anrecjob which schedules the execution of the playbook.
The implemented operator fulfills 8 out of the 10 requirements defined fully, with one being postponed due to time constraints, and one being decided against during development.
In conclusion the operator serves as a migration bridge but is not a full replacement for containerized workloads in Kubernetes.
